\documentclass{ximera}

\input{../../preamble.tex}

\title{Quadratic Forms}

\begin{document}

\begin{abstract}
complete the square
\end{abstract}
\maketitle




\subsection*{Quadratic Function Analysis}

What do we want to know when we analyze any function?

We want to know the 
\begin{itemize}
     \item \textbf{\textcolor{red!80!black}{Domain}} 
     \item \textbf{\textcolor{red!80!black}{Zeros}} 
     \item \textbf{\textcolor{red!80!black}{Continuity}} 
\begin{itemize}
     \item \textbf{\textcolor{purple!85!blue}{discontinuities}} 
     \item \textbf{\textcolor{purple!85!blue}{singularities}} 
\end{itemize}
     \item \textbf{\textcolor{red!80!black}{End-Behavior}} 
     \item \textbf{\textcolor{red!80!black}{Behavior}} 
\begin{itemize}
     \item \textbf{\textcolor{purple!85!blue}{intervals where increasing}} 
     \item \textbf{\textcolor{purple!85!blue}{intervals where decreasing}} 
\end{itemize}
     \item \textbf{\textcolor{red!80!black}{Global Maximum and Minimum}} 
     \item \textbf{\textcolor{red!80!black}{Local Maximums and Minimums}} 
     \item \textbf{\textcolor{red!80!black}{Range}} 
     \item \textbf{\textcolor{blue!55!black}{...and we would like a nice graph}} 
\end{itemize}



\textbf{\textcolor{blue!75!black}{Quadratic Domain and Range}} \\


The implied or natural domain of a quadratic function is all real numbers.  Of course, any quadratic function could come with a stated domain that restricts the domain. Or, a quadratic might be a piece of a piecewise defined function and only applied to a stated subset of the whole domain. Or, the quadratic function might be used as a model and the situation restricts the domain to an applied domain. \\


From the graphs above we can see that the natural range of a quadratic comes in two types.  

\begin{itemize}
\item The range could be all real numbers greater than or equal to some particular real number (the minimum).
\item The range could be all real numbers less than or equal to some particular real number (the maximum).
\end{itemize}

If there is a stated domain, then the range will be restricted accordingly. \\





\textbf{\textcolor{blue!75!black}{Quadratic Zeros}} \\

The number $0$ has a unique property appropriately named as the \textbf{Zero Product Property}.  



\begin{definition}  \textbf{\textcolor{green!50!black}{Zero Product Property}} \\

If $a$ and $b$ are real numbers and $a\cdot b = 0$, then either $a=0$ or $b=0$.

\end{definition}


Zero is the ONLY number with such an identifiable test, therefore, we use it...\textbf{\textcolor{purple!85!blue}{A LOT}}!   Like when solving equations. \\


\begin{procedure}  \textbf{\textcolor{blue!55!black}{Solving Equations}}

We have two methods for solving equations.

\begin{itemize}
     \item If you know what to do, then go do it. \\
     \item If you do not know what to do, then get everything to one side with $0$ on the other side and factor.
\end{itemize}


Unless we happen to know something special (which we often do), when solving equations our general approach is to set everything equal to zero and then convert the expression into a product.  This process is called \textbf{factoring}.  The whole idea is to make our situation look like the \textbf{zero product property}.

\end{procedure}






We have already seen that factors correspond to zeros of our function. Thus, it comes as no surprise that identifying zeros of functions is a top priority.  Factoring is a top priority. Spotting horizontal intercepts in the graph is a top priority.

As we can see from their parabolic graphs below, quadratics can have $0$, $1$, or $2$ real zeros.  How do we find them? Factoring is one way.  











\begin{center}
\textbf{\textcolor{green!50!black}{ooooo-=-=-=-ooOoo-=-=-=-ooooo}} \\

more examples can be found by following this link\\ \link[More Examples of Quadratics]{https://ximera.osu.edu/csccmathematics/precalculus1/precalculus1/projectileMotion/examples/exampleList}

\end{center}











\end{document}
